\section{Objetivos e Hipótesis}

\subsection{Hipótesis}

A partir del análisis de interacciones previas entre un usuario y un sistema informático, así como de las conversaciones sostenidas entre el usuario y un asistente virtual, es posible construir una base de datos híbrida — que combine componentes semánticos y estructurados — de la información personal del usuario. Dicha base de datos podría ser consultada posteriormente para aportar contexto relevante en la generación de respuestas por parte del asistente virtual.

El tema fundamental de la privacidad de los datos personales no es objeto de este trabajo ya que el foco está puesto en la funcionalidad y la precisión de la hiperpersonalización.

\subsection{Objetivos}

Se plantea como objetivo principal encontrar una solución tan simple, económica y eficiente como sea posible para la construcción de sistemas conversacionales que permitan la hiperpersonalización de las interacciones con los usuarios. Desde la perspectiva de las tecnoologías aplicadas, es preferible hallar una solución subóptima con las herramientas que se encuentran en uso actualmente que una solución óptima con tecnologías difíciles de integrar, agregando complejidad al mantenimiento de todo el sistema.